\documentclass[aspectratio=169, 11pt]{beamer}

% --- Theme ---
\usetheme{Madrid}
\usecolortheme{whale}

\definecolor{darkheader}{RGB}{44,62,80}
\definecolor{accentred}{RGB}{170,50,50}
\definecolor{accentgreen}{RGB}{39,118,72}
\definecolor{accentblue}{RGB}{30,80,150}
\setbeamercolor{frametitle}{bg=darkheader, fg=white}
\setbeamercolor{title}{bg=darkheader, fg=white}
\setbeamercolor{structure}{fg=darkheader}
\setbeamertemplate{navigation symbols}{}
\setbeamertemplate{footline}[frame number]

% --- Packages ---
\usepackage{amsmath, amssymb, amsthm}
\usepackage{booktabs, tcolorbox}
\usepackage{tikz}
\usetikzlibrary{arrows.meta, positioning}

% --- Macros ---
\newcommand{\E}{\mathbb{E}}
\newcommand{\F}{\mathcal{F}}
\newcommand{\cL}{\mathcal{L}}
\newcommand{\Dsigma}{\Delta\sigma_{IV}}
\newcommand{\phipi}{\phi_{\pi}}
\newcommand{\phiy}{\phi_{y}}
\newenvironment{citemize}{\begin{itemize}\setlength{\itemsep}{2pt}\setlength{\parskip}{0pt}}{\end{itemize}}

\newcommand{\contribbox}[1]{%
  \begin{tcolorbox}[colback=accentgreen!10,colframe=accentgreen,boxrule=0.8pt,arc=2pt,left=4pt,right=4pt,top=2pt,bottom=2pt]
    #1\end{tcolorbox}}

\newcommand{\riskbox}[1]{%
  \begin{tcolorbox}[colback=accentred!8,colframe=accentred,boxrule=0.8pt,arc=2pt,left=4pt,right=4pt,top=2pt,bottom=2pt]
    #1\end{tcolorbox}}

\newcommand{\fixbox}[1]{%
  \begin{tcolorbox}[colback=accentblue!8,colframe=accentblue,boxrule=0.8pt,arc=2pt,left=4pt,right=4pt,top=2pt,bottom=2pt]
    #1\end{tcolorbox}}

% --- Title ---
\title[Execution Plan v2]{Updated Execution Plan}
\subtitle{New contributions, execution risks, and mitigations}
\author{Aaron Huberman}
\institute{Duke University}
\date{\today}

% ============================================================
\begin{document}
% ============================================================

\begin{frame}
  \titlepage
\end{frame}

% ============================================================
\section{Project 1: Parameter Forward Guidance}
% ============================================================

\begin{frame}{Project 1 --- What Is New}
\small
\contribbox{\textbf{Headline.} Parameter FG as a structural resolution of the forward guidance puzzle. FOMC communication shifts a continuous posterior over $(\phipi,\phiy)$, not the rate level. Parameter-FG commitments are self-enforcing and generate appropriately sized yield responses where rate-level FG explodes.}

\vspace{0.2cm}
\textbf{Three new objects.}
\begin{citemize}
  \item \textbf{Three-way term premium.} Decompose into rate-path, $\phipi$, and $\phiy$ uncertainty components --- the latter two are absent from every existing term-structure model.
  \item \textbf{Divergence test.} Under parameter FG, a dovish communication compresses rate vol \emph{and} widens inflation breakeven vol simultaneously. Rate-level FG cannot generate this TIPS--SOFR divergence without an independent Fisher shock.
  \item \textbf{Asymmetric credibility.} The posterior-shift channel is incentive-compatible only during easing. Tightening episodes (2022--23) serve as a structural placebo: the channel should be off. Tested by pooling ECB/BoE threshold-guidance episodes to expand the sample beyond one US cycle.
\end{citemize}
\end{frame}

% ------------------------------------------------------------
\begin{frame}{Project 1 --- Risks, Prior Work, and Mitigations}
\footnotesize

\riskbox{\textbf{Risk 1.} Fisher critique: any rate-vol / breakeven-vol divergence could be mechanically generated by rate-level FG with a non-trivial Fisher equation.
\quad\textbf{Risk 2.} Tractability: continuous posteriors over $(\phipi,\phiy)$ have no standard fixed-point solution.
\quad\textbf{Risk 3.} NLP circularity: the communication signal $s_t$ may be endogenous to the very yield response it is meant to explain.}

\vspace{0.15cm}
\textbf{Why prior work does not solve this.}
\textbf{Bianchi \& Melosi (2017)} --- regime is binary, transition matrix constant; cannot accommodate a continuous, communication-driven posterior.
\textbf{Del Negro--Giannoni--Patterson (2023)} --- diagnoses the FG puzzle but stays within rate-level commitment; the parameter channel is not on the menu.
\textbf{Hansen, McMahon \& Tong (2019)} --- NLP on FOMC text to predict macro outcomes, not to extract a posterior over Taylor coefficients.
The macro-learning and asset-pricing literatures have never been forced to connect because the empirical signature (TIPS--SOFR divergence) was not available at scale until 2011--2020.

\vspace{0.15cm}
\fixbox{\textbf{Mitigations.}
\textbf{R1}: Calibrate an upper bound on Fisher-mechanical breakeven movement; reject rate-level FG where the observed divergence exceeds it.
\textbf{R2}: Restrict posteriors to Gaussian conjugate family; perturb around the mean belief so the equilibrium is solved once symbolically, not per event.
\textbf{R3}: Train the NLP extractor on textual features only (syntactic patterns around ``tolerate,'' ``even if,'' ``maximum employment''); validate against NY Fed Primary Dealer survey expectations, holding yields out of training entirely.}

\end{frame}

% ============================================================
\section{Project 2: Pricing Kernel with Normalizing Flows}
% ============================================================

\begin{frame}{Project 2 --- What Is New}
\small
\contribbox{\textbf{Headline.} Build the \emph{term structure} of pricing kernels $\hat{M}_t(R,\tau)$ at $\tau \in \{1,3,6,12\}$ months using horizon-specific conditional normalizing flows for $\hat{p}$. Test whether the puzzle is horizon-dependent and whether residual humps track horizon-matched belief dispersion.}

\vspace{0.2cm}
\textbf{Three new objects.}
\begin{citemize}
  \item \textbf{Conditional flow for $\hat{p}$.} Replace the unconditional historical distribution with $\hat{p}(R_{t+\tau}\mid\F_t)$ trained on VIX, credit spreads, realized variance, and term spread. Flows are trained jointly with $\tau$ as a conditioning input, pooling information across horizons.
  \item \textbf{Horizon-dependent puzzle.} $\hat{M}_t(R,\tau)$ at short horizons reflects rate-path disagreement; at long horizons, regime / inflation disagreement. The term structure of the hump pins down the economic source in a way a single-horizon estimate cannot.
  \item \textbf{Post-2020 VRP breakdown.} The VRP regression failure since 2020 is attributed to a structural shift in $\hat{p}_t$, not the kernel's shape --- the conditional flow provides the clean counterfactual $\hat{p}$ needed to run this diagnostic. Cross-section on individual stocks (regressing firm-level puzzle on analyst dispersion) validates that the channel is real and not an index artifact.
\end{citemize}
\end{frame}

% ------------------------------------------------------------
\begin{frame}{Project 2 --- Risks, Prior Work, and Mitigations}
\footnotesize

\riskbox{\textbf{Risk 1.} Ratio instability: small errors in $\hat{p}$ blow up in $q/\hat{p}$, especially in tails.
\quad\textbf{Risk 2.} $q$--$p$ alignment: $q_t$ is cross-sectional (options); $\hat{p}_t$ is time-series (returns). Consistent state space is non-trivial.
\quad\textbf{Risk 3.} Sample size: 300 monthly non-overlapping observations for the 12-month flow is thin for a deep generative model.
\quad\textbf{Risk 4.} Horizon-matched dispersion does not exist off-the-shelf.}

\vspace{0.1cm}
\textbf{Why prior work does not solve this.}
\textbf{Yang \& Hospedales (2023)}: $q$-side only; $\hat{p}$ and $M=q/p$ never formed.
\textbf{Jackwerth (2000); Rosenberg \& Engle (2002); A\"it-Sahalia \& Lo (2000)}: classical unconditional estimators, single horizon; per-horizon conditional $\hat{p}$ was statistically intractable.
\textbf{Christoffersen et al.\ (2010)}: parametric SV imposes functional form a flow learns automatically.
Neither the ML-finance crossover papers nor the PK literature has built a term structure of $\hat{M}$.

\vspace{0.1cm}
\fixbox{\textbf{Mitigations.}
\textbf{R1}: Work in log-density space ($\log q - \log \hat{p}$) and regularize tails with a Student-$t$ base distribution; bootstrap over training draws for confidence bands before claiming any hump.
\textbf{R2}: Align on $R = \log(S_{t+\tau}/S_t)$ with SPX options and CRSP VW returns under identical dividend treatment.
\textbf{R3}: Shared-weight architecture with $\tau$ as conditioning input pools information; only tail shape varies by horizon.
\textbf{R4}: Use $\gamma_t(\tau)$ from Project 3 --- the rotation-based curvature measure exists at any optionable maturity and is the primary cross-project integration point.}

\end{frame}

% ============================================================
\section{Project 3: Vol Smile Curvature}
% ============================================================

\begin{frame}{Project 3 --- What Is New}
\small
\contribbox{\textbf{Headline.} Decompose the change in the implied vol smile around macro announcements into $(\alpha_t,\beta_t,\gamma_t)$ at each maturity $\tau$. The maturity-profile of curvature $\gamma_t(\tau)$ is a structural, horizon-specific belief-dispersion instrument. A supervised classifier maps $(\alpha,\beta,\gamma)(\tau)$ to announcement type.}

\vspace{0.2cm}
\textbf{Three new objects.}
\begin{citemize}
  \item \textbf{Structural mapping.} Two-period, two-type CRRA heterogeneous-beliefs model gives a closed-form $(\alpha,\beta,\gamma) \leftrightarrow (\mu_d,\sigma_d^2)$ correspondence. Makes the rotation signature structural, not merely descriptive. RANK falsification: $\beta\approx\gamma\approx 0$ for pure rate shocks.
  \item \textbf{Classifier.} Trained on external labels (Bloomberg consensus dispersion, SPF changes, 50--100 economist-labeled gold-standard events). Tests whether the implied vol surface contains belief-heterogeneity information not visible in standard moments.
  \item \textbf{Term structure and pre-announcement.} $\gamma_t(\tau)$ separates rate-path (short) from regime/inflation (long) disagreement. Pre-announcement $\gamma$ reveals prior dispersion; $\Delta\gamma$ post-announcement measures the belief update. Both plug directly into Projects 1 and 2 as the dispersion instrument.
\end{citemize}
\end{frame}

% ------------------------------------------------------------
\begin{frame}{Project 3 --- Risks, Prior Work, and Mitigations}
\footnotesize

\riskbox{\textbf{R1.} Intraday data: EOD OptionMetrics is too noisy for narrow event windows; intraday is expensive and carries microstructure noise.
\quad\textbf{R2.} Long-maturity illiquidity: $\gamma_t(\tau)$ contaminated by stale quotes at $\tau\!\geq\!12$mo.
\quad\textbf{R3.} Classifier labels: self-labeling is circular; external proxies are noisy.
\quad\textbf{R4.} Closed-form structural mapping requires a two-type simplification referees will challenge.}

\vspace{0.1cm}
\textbf{Why prior work does not solve this.}
\textbf{Cont \& da Fonseca (2002)}: decomposes surface \emph{level}, not event-window \emph{change}; no beliefs interpretation.
\textbf{Bernanke \& Kuttner (2005); Lucca \& Moench (2015)}: level-only ($\alpha$), no rotation, no maturity profile.
\textbf{Kelly, P\'astor \& Veronesi (2016)}: implied-vol level around political events; one factor, no geometric decomposition.
Heterogeneous-beliefs models have never been asked to deliver \emph{surface-change} predictions; the combination of intraday options, a structural beliefs model, and post-2020 events has not appeared in one paper.

\vspace{0.1cm}
\fixbox{\textbf{Mitigations.}
\textbf{R1}: EOD OptionMetrics is the main paper; LiveVol intraday is robustness on FOMC/NFP/CPI only, filtered by volume and spread.
\textbf{R2}: Core grid $\tau\!\in\!\{1,3,6,12\}$mo (liquid SPX); 24/36mo are robustness only.
\textbf{R3}: External labels only (Bloomberg consensus $\sigma$, SPF changes); validate on 50--100 economist-labeled events.
\textbf{R4}: Two-type model is the minimal sufficient structure; sign of $(\alpha,\beta,\gamma)$ predictions survives three-type extension under weak regularity.}

\end{frame}

\end{document}
